\documentclass[final]{beamer}
%% Possible paper sizes: a0, a0b, a1, a2, a3, a4.
%% Possible orientations: portrait, landscape
%% Font sizes can be changed using the scale option.
\usepackage[size=a1,orientation=portrait,scale=1.1]{beamerposter}

\usetheme{gemini}
\usecolortheme{default}
\useinnertheme{rectangles}

\usepackage[english]{babel}
\usepackage[utf8]{inputenc}
%\usepackage{beamerthemesplit}
\usepackage{graphics,epsfig, subfigure}
\usepackage{url}
\usepackage{srcltx}
\usepackage{mathrsfs}
\usepackage{amsfonts}
\usepackage{amsmath}
\DeclareMathOperator\arctanh{arctanh}
\usepackage{amssymb}
\usepackage{bbm}
\usepackage{amsthm}
\usepackage{graphicx}
\usepackage{centernot}
\usepackage{caption}
\usepackage{braket}
\usepackage{lastpage}
\usepackage{setspace}
\usepackage{xcolor}
\usepackage{mathtools}
\usepackage{cancel}
\usepackage{booktabs}
\usepackage{tcolorbox}
\usepackage[normalem]{ulem}
\usepackage{anyfontsize}
\usepackage{t1enc}

\usepackage[square,sort,comma,numbers]{natbib}

\renewcommand{\ULdepth}{1pt}

\newcommand{\euler}[1]{\text{e}^{#1}}
\newcommand{\Real}{\text{Re}}
\newcommand{\Imag}{\text{Im}}
\newcommand{\supp}{\text{supp}}
\newcommand{\norm}[1]{\left\lVert #1 \right\rVert}
\newcommand{\abs}[1]{\left\lvert #1 \right\rvert}
\newcommand{\floor}[1]{\left\lfloor #1 \right\rfloor}
\newcommand{\Span}[1]{\text{span}\left(#1\right)}
\newcommand{\dom}[1]{\mathscr D\left(#1\right)}
\newcommand{\Ran}[1]{\text{Ran}\left(#1\right)}
\newcommand{\conv}[1]{\text{co}\left\{#1\right\}}
\newcommand{\Ext}[1]{\text{Ext}\left\{#1\right\}}
\newcommand{\vin}{\rotatebox[origin=c]{-90}{$\in$}}
\newcommand{\interior}[1]{%
	{\kern0pt#1}^{\mathrm{o}}%
}
\renewcommand{\braket}[1]{\left\langle#1\right\rangle}
\newcommand*\diff{\mathop{}\!\mathrm{d}}
\newcommand{\ie}{\emph{i.e.} }
\newcommand{\eg}{\emph{e.g.} }
\newcommand{\dd}{\partial }
\newcommand{\R}{\mathbb{R}}
\newcommand{\C}{\mathbb{C}}
\newcommand{\w}{\mathsf{w}}
\newcommand{\rr}{\mathcal{R}}


\newcommand{\Gliminf}{\Gamma\text{-}\liminf}
\newcommand{\Glimsup}{\Gamma\text{-}\limsup}
\newcommand{\Glim}{\Gamma\text{-}\lim}

\newtheorem{mtheorem}{Theorem}
\newtheorem{mdefinition}{Definition}
\newtheorem{mproposition}{Proposition}
\newtheorem{mlemma}{Lemma}
\newtheorem{mcorollary}{Corollary}
\newtheorem{mremark}{Remark}
\newtheorem{mconjecture}{Conjecture}


\definecolor{qmathblue}{RGB}{61,131,131}
\definecolor{kugreen1}{RGB}{50,93,61}
\definecolor{kugreen}{RGB}{70,116,60}
\definecolor{kugreenlys}{RGB}{132,158,139}
\definecolor{kugreenlyslys}{RGB}{173,190,177}
\definecolor{kugreenlyslyslys}{RGB}{239,249,240}
\definecolor{qmathbluelyslyslys}{RGB}{236,243,243}
\setbeamercovered{transparent}
\mode<presentation>
%\usetheme[numbers,totalnumber,compress,sidebarshades]{PaloAlto}
%\setbeamertemplate{footline}[frame number]
\setbeamertemplate{theorems}[numbered]

\usecolortheme[named=qmathblue]{structure}
\useinnertheme{circles}
\usefonttheme[onlymath]{serif}
\setbeamercovered{transparent}
\setbeamertemplate{blocks}[rounded][shadow=false]
\setbeamercolor{block body}{bg=qmathbluelyslyslys,fg=black}
\setbeamercolor{block title}{bg=qmathbluelyslyslys,fg=black}



\newlength{\sepwidth}
\newlength{\colwidth}
\setlength{\sepwidth}{0.03\paperwidth}
\setlength{\colwidth}{0.467\paperwidth}
\newcommand{\separatorcolumn}{\begin{column}{\sepwidth}\end{column}}

\newlength{\sepwidthh}
\newlength{\colwidthh}
\setlength{\sepwidthh}{0.006\paperwidth}
\newcommand{\separatorcolumnn}{\begin{column}{\sepwidthh}\end{column}}
%\useoutertheme{infolines} 
%% ====================
% BibLaTeX
% ====================

\usepackage[backend=biber,
	bibstyle=authoryear,
	citestyle=authoryear,
	style=authoryear,
	maxcitenames=2,
	maxbibnames=20, % limit the length of list of names (authors/editors/etc.)
	sorting=ydnt, % sort references by year (descending), name, title
	dashed=false, % show authors instead of dash in publications having the same authors
	giveninits=true % render authors' given name initials and not the full given names
]{biblatex}
%% Biblatex with Beamer bibliography icons
\setbeamertemplate{bibliography item}{%
	\ifboolexpr{ test {\ifentrytype{book}} or test {\ifentrytype{mvbook}}
		or test {\ifentrytype{collection}} or test {\ifentrytype{mvcollection}}
		or test {\ifentrytype{reference}} or test {\ifentrytype{mvreference}} }
	{\setbeamertemplate{bibliography item}[book]}
	{\ifentrytype{online}
		{\setbeamertemplate{bibliography item}[online]}
		{\setbeamertemplate{bibliography item}[article]}}%
	\usebeamertemplate{bibliography item}}
\defbibenvironment{bibliography}
{\list{}
	{\settowidth{\labelwidth}{\usebeamertemplate{bibliography item}}%
		\setlength{\leftmargin}{\labelwidth}%
		\setlength{\labelsep}{\biblabelsep}%
		\addtolength{\leftmargin}{\labelsep}%
		\setlength{\itemsep}{\bibitemsep}%
		\setlength{\parsep}{\bibparsep}}}
{\endlist}
{\item}
%% Redefine \refname
\renewcommand{\bibname}{References}
%% Redefine \parencite to use square brackets instead of braces
\DeclareCiteCommand{\parencite}
{\usebibmacro{prenote}}
{\usebibmacro{citeindex}%
	\printtext[bibhyperref]{[\usebibmacro{cite}]}}
{\multicitedelim}
{\usebibmacro{postnote}}
%% Highlight author names using Beamer data annotation
%% Usage: add a new line `author+an = {<author-order>=highlight}` to an entry
%% For example: author+an = {3=highlight} => highlight the 3rd author name
\AtBeginBibliography{
	\renewcommand*{\mkbibnamegiven}[1]{%
		\ifitemannotation{highlight}
		{\textbf{#1}}
		{#1}%
	}
	
	\renewcommand*{\mkbibnamefamily}[1]{%
		\ifitemannotation{highlight}
		{\textbf{#1}}
		{#1}%
	}
}

% ====================
% Boxes with tcolorbox
% ====================
\usepackage[most]{tcolorbox}

%%% Beamer colors in boxes

\newcommand{\beamercolorsinboxes}[1]{
	\setbeamercolor{itemize item}{fg=#1!75!black}
	\setbeamercolor{itemize/enumerate body}{fg=#1!65!white}
	\setbeamercolor{itemize/enumerate subbody}{fg=#1!65!white}
	\setbeamercolor{item projected}{fg=white, bg=#1!75!black}
}

%%% Highlight Oval Box
\newtcbox{\xmybox}[1][red]{on line,
	arc=7pt,colback=#1!10!white,colframe=#1!50!black,
	before upper={\rule[-3pt]{0pt}{10pt}},boxrule=1pt,
	boxsep=0pt,left=6pt,right=6pt,top=2pt,bottom=2pt}
%%% Box for stating problems
%%%%%%%%
%Usage: (similar for infobox)
%	\begin{defbox}{title}
	%		contents
	%	\end{defbox}
%%%%%%%%
\newtcolorbox{defbox}[1]{%
	enhanced,
	attach boxed title to top 	left={xshift=5mm,yshift=-5mm,yshifttext=-5mm},
	colback=cyan!5!white,
	colframe=cyan!75!black,
	coltitle=cyan!80!black,
%	left=0mm,right=0mm,top=2mm,bottom=0mm,
	title={#1},
	fonttitle=\bfseries\large, fontupper=\color{cyan!65!white},
	boxed title style={colback=cyan!5!white,colframe=cyan!75!black},
	before upper={
		\beamercolorsinboxes{cyan}
	}
}%
%%% Box for announcement
\newtcolorbox{infobox}[1]{%
	enhanced,
	attach boxed title to top 	left={xshift=5mm,yshift=-5mm,yshifttext=-5mm},
	colback=yellow,
	colframe=red!75!black,
	coltitle=red!75!black,
%	left=0mm,right=0mm,top=2mm,bottom=0mm,
	title={#1},
	fonttitle=\bfseries\large, fontupper=\color{red!65!white},
	boxed title style={colback=yellow,colframe=red!75!black},
	before upper={
		\beamercolorsinboxes{red}
	}
}%
%%% Box for example
\newtcolorbox{exabox}[1]{%
	enhanced,
	attach boxed title to top 	left={xshift=5mm,yshift=-5mm,yshifttext=-5mm},
	colframe=brown!75!black,colback=brown!5!white,coltitle=brown!50!brown!75!black,
%	left=0mm,right=0mm,top=2mm,bottom=0mm,
	title={#1},
	fonttitle=\bfseries\large, fontupper=\color{brown!65!white},
	boxed title style={colback=brown!5!white,coltitle=brown!50!brown!75!black},
	before upper={
		\beamercolorsinboxes{brown}
	}
}%
%%% Theorem Box
%%%%%%%%
%Usage: (similar for conjecture, lemma, etc.)
%	\begin{thm}{title}{nameref}
	%		contents
	%	\end{thm}
% Use \ref{thm:nameref} to refer to the theorem
%%%%%%%%
%%%% Use \newtcbtheorem[number within=section]{thm} to number within each section
\newtcbtheorem[]{thm}%
{Theorem}{attach boxed title to top 	left={xshift=5mm,yshift=-5mm,yshifttext=-5mm},
	enhanced jigsaw,
	%	top=2mm,bottom=0mm,left=0mm,right=0mm,
	fonttitle=\bfseries\large,fontupper=\itshape\color{blue!65!white},
	colframe=blue!75!black,colback=blue!5!white,coltitle=blue!50!blue!75!black,
	boxed title style={colback=blue!5!white,coltitle=blue!50!blue!75!black},
	before upper={
		\beamercolorsinboxes{blue}
	}
}{thm}%
%%% Proposition Box
\newtcbtheorem[use counter from=thm]{prop}%
{Proposition}{attach boxed title to top 	left={xshift=5mm,yshift=-5mm,yshifttext=-5mm},
	enhanced jigsaw,
	%	top=2mm,bottom=0mm,left=0mm,right=0mm,
	fonttitle=\bfseries\large,fontupper=\itshape,
	colframe=gray!75!black,colback=gray!5!white,coltitle=gray!50!gray!75!black,
	boxed title style={colback=gray!5!white,coltitle=gray!50!gray!75!black},
	before upper={
		\beamercolorsinboxes{gray}
	}
}{prop}%
%%% Conjecture Box
\newtcbtheorem[use counter from=thm]{conj}%
{Conjecture}{attach boxed title to top 	left={xshift=5mm,yshift=-5mm,yshifttext=-5mm},
	enhanced jigsaw,
	%	top=2mm,bottom=0mm,left=0mm,right=0mm,
	fonttitle=\bfseries\large,fontupper=\slshape,
	colframe=orange!75!black,colback=orange!5!white,coltitle=orange!50!orange!75!black,
	boxed title style={colback=orange!5!white,coltitle=orange!50!orange!75!black},
	before upper={
		\beamercolorsinboxes{orange}
	}
}{conj}%
%%% Lemma Box
\newtcbtheorem[use counter from=thm]{lem}%
{Lemma}{attach boxed title to top 	left={xshift=5mm,yshift=-5mm,yshifttext=-5mm},
	enhanced jigsaw,
	%	top=2mm,bottom=0mm,left=0mm,right=0mm,
	fonttitle=\bfseries\large,fontupper=\itshape,
	colframe=green!75!black,colback=green!5!white,coltitle=green!50!green!75!black,
	boxed title style={colback=green!5!white,coltitle=green!50!green!75!black},
	before upper={
		\beamercolorsinboxes{green}
	}
}{lem}%
%%% Claim Box
\newtcbtheorem[use counter from=thm]{clm}%
{Claim}{attach boxed title to top 	left={xshift=5mm,yshift=-5mm,yshifttext=-5mm},
	enhanced jigsaw,
	%	top=2mm,bottom=0mm,left=0mm,right=0mm,
	fonttitle=\bfseries\large,fontupper=\itshape,
	colframe=pink!75!black,colback=pink!5!white,coltitle=pink!50!pink!75!black,
	boxed title style={colback=pink!5!white,coltitle=pink!50!pink!75!black},
	before upper={
		\beamercolorsinboxes{pink}
	}
}{clm}%

\logoleft{\hspace*{2cm}\includegraphics[width=10cm]{qmath.jpg}}
\logoright{\hspace*{-5cm}\includegraphics[width=5cm]{kuscience-logo}}

\title{The ground state energy of 1D dilute many-body quantum systems}

\author{\underline{Johannes Agerskov}\textsuperscript{1}, Robin Reuvers\textsuperscript{2}, and Jan Philip Solovej\textsuperscript{1}}
\institute{1. Department of Mathematics, University of Copenhagen, Universitetsparken 5, DK-2100 Copenhagen \O, Denmark
	\\
	2. Universit\`{a} degli Studi Roma Tre, Dipartimento di Matematica e Fisica, L.go S. L. Murialdo 1, 00146 Roma, Italy}
\date{June 15th 2022}

\setbeamercovered{invisible}

\begin{document}
	\begin{frame}[t]
		\vspace*{-1.5cm}
		\begin{columns}[t]
			\begin{column}{2\colwidth+\sepwidthh}
				\begin{tcolorbox}[colframe=qmathblue,colback=qmathbluelyslyslys,title=\centering Abstract]
					We study the ground state energy of a gas of 1D bosons with density $\rho$, interacting through a general, repulsive 2-body potential with scattering length $a$, in the dilute limit $\rho |a|\ll1$. The first terms in the expansion of the thermodynamic energy density are $\pi^2\rho^3/3(1+2\rho a)$, where the leading order is the 1D free Fermi gas. This result covers the Tonks--Girardeau limit of the Lieb--Liniger model as a special case, but given the possibility that $a>0$, it also applies to potentials that differ significantly from a delta function.
				\end{tcolorbox}
			\end{column}
			
		\end{columns}
		
	\begin{columns}[t]
		\separatorcolumn
		
		\begin{column}{\colwidth}
		\begin{tcolorbox}[colframe=qmathblue,colback=qmathbluelyslyslys,title=Set up and previous results]
		We consider the model \begin{equation}
			H=-\sum^N_{i=1}\Delta_{x_i}+\sum_{1\leq i<j\leq N}v(x_i-x_j),
		\end{equation}	
		where $ \Delta_{x_i} $ is the $ d $-dimension Laplacian w.r.t $ x_i\in\R^d $, and $ v $ is a repulsive, spherically symmetry potential.\\
		The corresponding bosonic energy functional is 
		\begin{equation}
		\mathcal{E}(\psi)=\int_{\Lambda_L}\left(\sum_{i=1}^{N}\abs{\nabla_i\psi}^2+\sum_{i<j} v_{ij}\abs{\psi}^2\right)\quad \text{on } L^2(\Lambda_L)^{\otimes_{\text{sym}} N},
		\end{equation}
		with $ v_{ij}=v\left(\abs{x_i-x_j}\right) $, and the ground state energy is defined by 
		\begin{equation}
			E(N,L)\coloneqq\inf_{\psi\in\mathcal{D}(\mathcal{E}),\ \norm{\psi}^2=1}\mathcal{E}(\psi).
		\end{equation}
		Let $ e(\rho)\coloneqq\lim\limits_{\substack{L\to\infty\\ N/L^{d}\to\rho}}E(N,L)/L^{d} $, then previously known result in $ d=2,3 $ are 
		\begin{theorem}[$ d=3 $ result, Lee-Huang-Yang \cite{lee1957eigenvalues}, \cite{yau2009second,fournais2020energy,basti2021new,fournais2021energy}]
			\begin{equation}
			e(\rho)=4\pi\rho^2 a\left(1+\frac{128}{15\sqrt{\pi}}\sqrt{\rho a^3}+o(\sqrt{\rho a^3})\right).
			\end{equation}
		\end{theorem}
		
		\begin{theorem}[$ d=2 $ result, \cite{lieb2001ground}]
			\begin{equation}
			e(\rho)=4\pi \rho^2\left(\abs{\ln(\rho a^2)}^{-1}+o(\abs{\ln(\rho a^2)}^{-1})\right) .
			\end{equation}
		\end{theorem}
		
		
		\end{tcolorbox}
		
		\begin{tcolorbox}[colframe=qmathblue,colback=qmathbluelyslyslys,title=The scattering length]
			\begin{theorem}[\cite{lieb2001ground}]
				For $ B_R\coloneqq\{0<\abs{x}<R\}\subset \R^d $ with $ R>R_0\coloneqq\text{range}(v) $, let $ \phi\in H^1(B_{R}) $ satisfy
				\begin{equation}
				-\Delta \phi +\frac12 v\phi=0,\qquad \text{on }B_R,
				\end{equation}
				with boundary condition $ \phi(x)=1 $ for $ \abs{x}=R$.
				Then $ \phi(x)=f(\abs{x}) $ for some $ f:(0,R]\to [0,\infty) $, and for $ \text{range}(v)<r<R $, we have \begin{equation}
				f(r)=\begin{cases}
				(r-a)/(R-a) &\text{for }d=1\\
				\ln(r/a)/\ln(R/a) &\text{for }d=2\\
				(1-ar^{2-d})/(1-aR^{2-d})&\text{for }d\geq 3,
				\end{cases}
				\end{equation}
				with some constant $ a $ called the (s-wave, or even-wave in $ d=1 $) \textbf{scattering length}.
			\end{theorem}
			Similarly, one can define the p-wave, or odd-wave in $ d=1 $, scattering length by having a p-wave boundary condition at $ \abs{x}=1 $.
		\end{tcolorbox}
			
		\begin{tcolorbox}[colframe=qmathblue,colback=qmathbluelyslyslys,title=Main result]
			Let $ v=v_{\text{reg}}+v_{\text{h.c.}} $, with $ v_{\text{reg}} $ a finite measure, and $ v_{\text{h.c.}} $ a positive linear combination of hard core potentials of the form 
			\begin{equation*}
			v_{[x_1,x_2]}(x):=\begin{cases*}
			\ \infty& \phantom{when} $|x|\in [x_1,x_2]$\\
			\ 0 & \phantom{when} \text{otherwise}
			\end{cases*}.
			\end{equation*}
			Then we have the result
			\begin{theorem}[Bosons {[A., R. Reuvers, J. P. Solovej, 2022]}]
				\label{TheoremMain}
				Consider a one dimensional ($ d=1 $) Bose gas with repulsive interaction  $v=v_{\text{reg}}+v_{\text{h.c.}}$ as defined above. Write $\rho=N/L$. For $\rho|a|$ and $\rho R_0$ sufficiently small, the ground state energy can be expanded as 
				\begin{equation}
				\label{result}
				E(N,L)=N\frac{\pi^2}{3}\rho^2\left(1+2\rho a+
				\mathcal{O}
				\left((\rho|a|)^{6/5}+(\rho R_0)^{6/5}+N^{-2/3}\right)\right),
				\end{equation}
				where $a$ is the scattering length of $v$.
			\end{theorem}
			Since spinless (or spin-aligned) fermions are unitarily equivalent to impenetrable bosons, we have the following theorem as a corolarry 
			\begin{theorem}[Spinless fermions {[A., R. Reuvers, J. P. Solovej, 2022]}]\label{TheoremFermion}
				Consider a spinless Fermi gas with repulsive interaction  $v=v_{\text{reg}}+v_{\text{h.c.}}$ as defined before. Let $ E_F(N,L)$ be its associated ground state energy. Write $\rho=N/L$. For $\rho a_o$ and $\rho R_0$ sufficiently small, the ground state energy can be expanded as 
				\begin{equation}
				E_F(N,L)=N\frac{\pi^2}{3}\rho^2\left(1+2\rho a_o+\mathcal{O}\left(\left(\rho R_0\right)^{6/5}+N^{-2/3}\right)\right),
				\end{equation}
				where $ a_o\geq0 $ is the odd wave scattering length of $v$. 
			\end{theorem} 
		\end{tcolorbox}	
		
		\begin{tcolorbox}[colframe=qmathblue,colback=qmathbluelyslyslys,title=Upper bound proof sketch]
			Use Variational principle
			$$
			E(N,L)\leq \frac{\mathcal{E}(\Psi)}{\norm{\Psi}^2},\quad \text{for any }  \Psi\in \mathcal{D}(\mathcal{E}) .
			$$
			Let $ M=N(\rho b)^{3/2} $ and choose the trial state $ \Psi=\prod_{i=1}^{M}\Psi_i $, where \begin{equation}
				\Psi_i(x)=\begin{cases}
				\omega(\rr(x))\frac{\abs{\Psi^i_F(x)}}{\rr(x)}& \text{if }\rr(x)<b\\
				\abs{\Psi^i_F(x)}&\text{if }\rr(x)\geq b,
				\end{cases}
			\end{equation}
			and $ \omega $ is the (suitably normalized) solution to the two-body scattering equation,  $\Psi^i_F$ is the free Fermi ground state with $ N/M $ particles in the box $ [(i-1)L/M,iL/M] $, and $ \rr(x)\coloneqq \min_{i<j}(\abs{x_i-x_j}) $ is uniquely defined a.e. Trial state energy computation amounts, after some manipulations and integration by parts, to computing integrals involving reduced densities and reduced density matrices of the free Fermi gas.
		\end{tcolorbox}
		\end{column}
	\separatorcolumnn

	\begin{column}{\colwidth}
		\begin{tcolorbox}[colframe=qmathblue,colback=qmathbluelyslyslys,title=Upper bound proof sketch]
			For the free Fermi gas, then one-particle reduced density matrix is known
			\begin{equation}
			\begin{aligned}
			\gamma^{(1)}(x,y)&=\frac{2}{l}\sum_{j=1}^{N}\sin\left(\frac{\pi}{l}jx\right)\sin\left(\frac{\pi}{l} jy\right)\\
			&=\frac{\pi}{l}\left(D_{N}\left(\pi\frac{x-y}{l}\right)+D_{N}\left(\pi\frac{x+y}{l}\right)\right),
			\end{aligned}
			\end{equation}
			where $ l=L/M $ and $ D_N(x)=\frac{1}{2\pi}\sum_{k=-N}^{N}\euler{ikx}=\frac{\sin((N+1/2)x)}{2\pi\sin(x/2)} $ is the Dirichlet kernel.
			Bounds on $ k $-particle reduced density matrices are provided by Wick's theorem, \eg:
			\begin{mlemma}\label{Lemma rho2 bound}
				$$ \rho^{(2)}(x_1,x_2)\leq\left(\frac{\pi^2}{3}\rho^4+f(x_2)\right)(x_1-x_2)^2+\mathcal{O}(\rho^6(x_1-x_2)^4), $$ 
				with $ \int f(x_2)\diff x_2\leq \textnormal{ const. }\rho^3\log(N). $
			\end{mlemma}
			Collecting everything and choosing $ b=\max(\rho^{-1/5}\abs{a}^{4/5}, R_0) $ we find
			
			\begin{mproposition}[Upper bound Theorem \ref{TheoremMain}]
				\label{PropositionUpperBound}
				There exists a constant $C_\text{U}>0$ such that for $\rho|a|$, $\rho R_0\leq C_U^{-1}$, the ground state energy $E^D(N,L)$ satisfies
				\small{\begin{equation}
					\label{equpper}
					E^D(N,L)\leq N\frac{\pi^2}{3}\rho^2\left(1+2\rho a + C_\text{U}\left((\rho\abs{a})^{6/5}+(\rho R_0)^{3/2}+N^{-1}\right)\right).
					\end{equation}}
			\end{mproposition}
		\end{tcolorbox}
		
	\begin{tcolorbox}[colframe=qmathblue,colback=qmathbluelyslyslys,title=Lower bound proof sketch]
		Proof of lower bound consists of the following steps:
		\begin{itemize}
			\item Use Dyson's lemma to reduce to a nearest neighbor double delta-barrier potential.
			\item Reduce to the Lieb Liniger model by discarding \textbf{a small part} of the wave function.
			\item Use a known lower bound for the Lieb Liniger model.
		\end{itemize}	
		\begin{mlemma}[Dyson]\label{LemmaDyson} Let $ R>R_0=\textnormal{range}(v) $ and $ \varphi\in H^1(\R) $, then for any interval $ \mathcal{I}\ni 0 $ 
			\begin{equation}
			\int_{\mathcal{I}} \abs{\partial \varphi}^2+\frac12 v\abs{\varphi}^2\geq \int_{\mathcal{I}}\frac{1}{R-a}\left(\delta_R+\delta_{-R}\right)\abs{\varphi}^2,
			\end{equation}
			where $ a $ is the s-wave scattering length.
		\end{mlemma}
		Define $ \psi\in L^2([0,\ell-(n-1)R]^n) $ by 
				$$ \psi(x_1,x_2,...,x_n)=\Psi(x_1,R+x_2,...,(n-1)R+x_n), $$
				for $ x_1\leq x_2\leq...\leq x_n $ and symmetrically extended.\\\vspace{0.2cm}
				Then by Lemma \ref{LemmaDyson} \begin{equation}
				\begin{aligned}
				\mathcal{E}(\Psi)&\geq E^N_{LL}(n,\ell-(n-1)R,2/(R-a))\braket{\psi|\psi}\\
				&\geq n\frac{\pi^2}{3}\rho^2\left(1+2\rho(a-\cancel{R})+\cancel{2\rho R}-\text{const. }\frac{1}{N^{2/3}}\right)\braket{\psi|\psi}.
				\end{aligned}
				\end{equation}
		Mass of $ \psi $ can be bounded by 
		\begin{mlemma}\label{LemmaImprovedMassBound}
			For some constant $C>0$ and for $ n(\rho R)^2\leq  \frac{3}{16\pi^2}C $, $ \rho R\ll 1 $ and $ R>2\abs{a} $ we have
			\begin{equation}\label{EqImprovedMassBound}
			\begin{aligned}
			\braket{\psi|\psi} \geq 1-\textnormal{const. }\left(n(\rho R)^3+n^{1/3}(\rho R)^2\right).
			\end{aligned}
			\end{equation}
		\end{mlemma}
		and we find \begin{mproposition}\label{PropositionLowerBoundSpecN}
			For assumptions as in lemma \ref{LemmaImprovedMassBound} we have \begin{equation}
			E^N(n,\ell)\geq n\frac{\pi^2}{3}\rho^2\left(1+2\rho a+\textnormal{const. }\left(\frac{1}{n^{2/3}}+n(\rho R)^3+n^{1/3}(\rho R)^2\right)\right).
			\end{equation}
		\end{mproposition}
		Error grows with $n$, so localization is needed. Localization is justified by going to the grand canonical ensemble. We localize to boxes with $n=(\rho R)^{-9/5}$ giving
		\begin{mproposition}[Lower bound Theorem \ref{TheoremMain}]
			\label{PropositionLowerBound}
			There exists a constant $C_\text{L}>0$ such that the ground state energy $E^N(N,L)$ satisfies
			\begin{equation}
			\label{eqlower}
			E^N(N,L)\geq N\frac{\pi^2}{3}\rho^2\left(1+2\rho a-C_\text{L}\left((\rho\abs{a})^{6/5}+(\rho R_0)^{6/5}+N^{-2/3}\right)\right).
			\end{equation}
		\end{mproposition}
	\end{tcolorbox}
	\begin{tcolorbox}[colframe=qmathblue,colback=qmathbluelyslyslys,title=Spin-1/2 conjecture]
		For the spin-1/2 case, there are two solvable models. The Yang-Gaudin model, and the hard core model. Based on the ground state energies of these, one migth conjecture
		\begin{mconjecture}
		For assumptions as in Theorem \ref{TheoremMain}, the spin-1/2 Fermi ground state energy can be expanded as
			\begin{equation}
			\begin{aligned}
			E(N,L)=N\frac{\pi^2}{3}\rho^2\Big(1+2\ln(2)\rho a_e+2(1-\ln(2))\rho a_o\\+\mathcal{O}\left((\rho\max(\abs{a_e},a_o))^2\right)\Big).
			\end{aligned}
			\end{equation}
		\end{mconjecture}
	\end{tcolorbox}
	\begin{tcolorbox}[colframe=qmathblue,colback=qmathbluelyslyslys,title=References]
		\vspace*{0.06 cm}
		\bibliographystyle{amsplain}
		{\fontsize{12}{14.1}\selectfont\bibliography{bibliography}}
		\vspace*{0.06 cm}
	\end{tcolorbox}
	\end{column}	
	\separatorcolumn
	\end{columns}
	\end{frame}
\end{document}